%\input ../util/bookmacros
%\input ../util/docparams

\appendix{Appendix A}{Functions, Formulas, Graphs, and Notation}

In the main text we freely wrote formulas like $f(x) = x^2$ and drew
graphs of functions -- so you may well ask: are the formula, the graph,
and the function all the same thing? For the patchwork functions and
sequences of this book the distinction matters, and this appendix pins
it down. As
mentioned earlier in this book functions are rather abstract things. Functions
are not numbers or a collection of numbers; they are rules which assign one 
number to another. As such, we are tempted to identify anything tangibly 
associated with a function {\it as the function\/}. The two things that 
people often confuse with a function are formulas 
-- usually used to define a function -- and the graph of a function. 
There is a third problem which comes from the way mathematicians and others 
notate functions. 

Before we clarify these concepts let us emphasize that 
although functions are not formulas, the kinds of functions which make 
calculus computationally tractable are the ones 
which are defined by formulas or are made from a patch work of formulas.

\section{Functions are NOT Formulas}
You will often hear the terms {\it function\/} and {\it formula\/} used 
interchangeably. One may hear 
``This is a function that describes such and such.'' Or, 
``Here is a formula for computing such and such.'' Many people think 
they are the same thing. In fact, mathematicians a few centuries back 
{\it did\/} use 
the two terms interchangeably. They were driven to separate these notions 
through the problems they faced over the last few centuries --
just as we are being driven to introduce the concepts of Calculus by a 
particular set of problems. A few definitions will help draw this distinction.

\proclaim Definition(Function). A function, $f$, takes as input a number
from a set called the {\it domain\/} of $f$, and produces
{\it one\/} number as output. The set of outputs of $f$ (using input values from the domain) is called the 
{\it range\/} of $f$.

A simple example is the function which squares its input. If we call this
function $f$, then $f$ is defined by $f(x) = x^2$. We can take the domain of 
this function to be any set where the squaring operation makes sense. One usually
takes the domain to be the largest possible set one can. In this case, the 
domain is usually taken to be all real numbers.
As a second example, consider 
the function which takes the square root of its input. If we call this function
$g$, then $g$ is defined by $g(x) = \sqrt x$. We can take the domain of 
this function to be any set where the square-root operation makes sense. Since 
taking square roots of a number is only valid for non-negative numbers, we 
can take the domain of this function to be any subset of the non-negative numbers.
One usually takes the domain to be the non-negative numbers.
Lastly, consider the family of functions, $f(x) = x^{n}$ with integer $n$, $n < 0$.
The domain of these functions is taken to be the set of all numbers except $0$; that is,
all places where the function makes sense. 

A formula doesn't have much of a formal definition, but we can define it 
roughly as:

\proclaim Definition(Formula). A formula is an ``algebraic'' 
implementation of a 
function. That is, a formula is an ``algebraic rule'' which takes a numeric 
input and produces a numeric output.

Examples of formulas are:

\beginEnum
\enum{$1 + 2x$}
\enum{$x^2 + 3x + 2$}
\enum{$\sin^2(x) + x + \tan(x)$}
\enum{${1 + x^2 \over 2x - 7} e^x + 2 \tan(x)$}
\endEnum

These formulas can be used to define functions. Listed below are the associated 
functions from the formulas above.

\beginEnum
\enum{$f_1(x) = 1 + 2x$}
\enum{$f_2(x) = x^2 + 3x + 2$}
\enum{$f_3(x) = \sin^2(x) + x + \tan(x)$}
\enum{$f_4(x) = {1 + x^2 \over 2x - 7} e^x + 2 \tan(x)$}
\endEnum

It may seem like we're quibbling over trivial matters. To drive home the point 
that they are different we present below functions whose definitions go
beyond a single formula in various ways -- patchworks of formulas,
restricted domains, or no formula at all. In some examples,
a graph picturing the function is given.

{\bf Example 1.\/} $f$ is the absolute value function defined on the real 
line. There is a simple formula for this function. $f(x) = |x|$. But it 
may be defined as a compilation of even simpler formulas as:
$$
f(x) = \cases{ x & if $ x \ge 0$; \cr
-x & otherwise. \cr}
$$
\epsfigure{eps/absolute.eps}{The absolute value function $f(x) = |x|$.}

{\bf Example 2.\/} $g$ is referred to 
as a piecewise linear function. It is defined on the real line as:
$$
g(x) = \cases{ 1 & if $x < 0$; \cr
1+x & if $0 \le x < 3$; \cr
4 + (3 - x) & if $3 \le x < 5$; \cr
2 & if $x \ge 5$. \cr}
$$
\epsfigure{eps/piecewise.eps}{A sample piecewise linear function.}

{\bf Example 3.\/} $h$ is defined only on the interval $[0,10]$. That is, its domain is the set $[0,10]$.
This is an example of a discontinuous function; it makes a sudden jump when $x = 5$.
$$
h(x) = \cases{ x & if $ 0 \le x < 5$; \cr
7 & $5 \le x \le 10$. \cr}
$$
\epsfigure{eps/discount.eps}{A discontinuous function.}

{\bf Example 4.\/} $d$ is defined only on the input set $\{1,2,3\}$ (its domain) 
producing values in the output set $\{4,10,21\}$ (its range). 
$$
d(x) = \cases{ 10 & if $ x = 1$; \cr
21 & if $x = 2$; \cr
4 & if $x = 3$.}
$$

{\bf Example 5.\/} $f$ is defined only on the set of positive integers 
producing the even numbers.
$$
f(n) = 2\, n 
$$

{\bf Example 6.\/} $f$ is defined only on the set of positive integers 
producing the prime numbers.
$$
f(n) = p_n \quad {\rm where}\; p_n\; {\rm is} \;{\rm the} \;n^{\rm th}\; {\rm prime} 
$$
Several of the functions above, while not given by a single formula, are
patchworks of formulas (Examples 3 and 4). Example 5 is given by a simple
formula, $f(n) = 2n$, but on a restricted domain. Example 6 is a function for which one can't write down
a formula; nor a (finite) patchwork of formulas as there is no known formula for
expressing the sequence of prime numbers.%
\footnote{\kern 0.5pt \raise 0.5ex \hbox{\dag}}{Actually, a relatively concise formula was produced in the 1960s but it is not a practical means of computing primes.}

Let us also try to remove the confusion about functions and their 
descriptions. The function $f$ defined by $f(x) = x^2$ can also be 
defined by $f(z) = z^2$. Why? Because both descriptions tell us to do the 
same thing to their inputs. 
That is, the input variable is a dummy variable, not unlike a dummy variable 
used in a ``for loop'' of a programming language. For instance, consider 
the two ``for loops'' below: 

{\tt sum = 0; for(i = 0; i < 10; i++) $\{$ sum += i; $\}$}

{\tt sum = 0; for(j = 0; j < 10; j++) $\{$ sum += j; $\}$}

The results don't depend on $i$ or $j$. In the same way, be careful not to 
get too attached to the dummy variables used in function definitions.
We state this one more time with emphasis:

{\bf NOTE:\/} The following all define a function named $f$ that squares its 
input:
$$
f(x) = x^2 \;,\quad f(z) = z^2 \;,\quad f(a) = a^2 \;,\quad f(c) = c^2 
$$

\section{Functions are NOT Graphs}
Functions are abstract: a function is not a number but a rule that takes in one number and produces another;
it is difficult to get a picture of this. Graphs of functions give us a picture. However, 
it is worth mentioning again that graphs are NOT functions. A function 
has a graph, but a graph may not correspond to a function. To make 
the distinction clearer we need to define what we mean by a graph.

\proclaim Definition(Graph). A graph is a set of pairs of real numbers.
That is, a graph is a {\it subset\/} of $\{\,(x,y)\mid x,y\in~R\,\}$.\footnote{\kern 0.5pt \raise 0.5ex \hbox{\ddag}}%
{The symbol $R$ stands for the set of all real numbers.}

\proclaim Definition(Graph of a Function). The graph of a function, $f$, 
is the set $\{\,(x,f(x)) \mid x \in~{\rm Domain}(f)\,\}$.

This leads to a question: Given a graph, $G$, is there a corresponding 
function, $f$, for which the graph of $f$ is $G$?

The answer, in general, is no. Consider, for example, a graph that 
is a circle. 
The graph of a function has the property that for any 
given input, $x$, (provided it's in the function's domain) only one output is produced. 
What does this mean for the graph
of a function? It means that if we take a vertical line through $x$, it 
should only hit the graph of the function at one point. Because otherwise, 
if it hit the graph at more than one point, there would be more than
one value produced by $f$ -- and this can't be. This criterion is called
the {\it \index{vertical line test}\/}: a graph is the graph of a function
exactly when every vertical line hits it at most once. So, in the case of the circle
graph, we know that this is not the graph of a function because a vertical 
line hits more than one point for at least one value of $x$. In fact, a 
vertical line hits two points of the circle at all but the bounding $x$ values.

\section{Function Notation}
In many respects mathematics is a triumph of notation. We use symbols to name 
and manipulate abstract ideas. It is often not as precise as mathematicians 
would like, however. There are cases where the short 
hand of notation is confusing 
or misleading. One case in particular is the way people refer to 
functions.
In this book you will see functions notated as variables are notated; for example, we might call a function $f$.
 In some texts there is a concern that this simple name may confuse a function name with the name 
of a variable that has a numeric value. In order to add information 
(in this case to indicate that $f$ is a function)  some 
people use the notation $f(x)$ to try to make things clear. The sense which is 
to be conveyed by this notation is that $f$ is a function because
it takes inputs. The symbol $x$ is often used because it conveys the 
idea of ``variable input''; that is, $f$ can be fed any number as input.
This leads to confusion because the usual interpretation of $f(x)$ is $f$ evaluated at $x$.%
\footnote{\kern 0.5pt \raise 0.5ex \hbox{\dag}}{The only place where $f(x)$ has the sense of $f$ being fed an arbitrary input is when
we {\it define\/} the function $f$ as in $f(x) = $ value, or expression involving $x$. }
Other people will use the notation $f(\cdot)$ to more clearly indicate that $f$ is a function.
The ``dot'' serves as a place holder indicating that $f$ feeds on something; that is, it is a function.

Reiterating, we will use the notation $f(x)$ to 
mean: evaluate the function $f$ at $x$.

